% English counterpart of draft/background_and_motivation.md.
% Compile from paper/; see motivation-preview.tex for EuroSys 2027 settings.
\section{Background and Motivation}
\label{sec:background-motivation}

Multi-agent applications complete a user task through interdependent LLM
requests. After a worker failure, rebuilding lost key--value (KV) state can
restore an individual request, but the user-visible penalty depends on when
the entire workflow completes. Workflow dependencies and the cost of
rebuilding long contexts motivate allocating limited KV protection resources
to reduce workflow completion time (WCT).

\input{sections/motivation-figures}

\subsection{Workflows, Failures, and KV Recovery}
\label{sec:motivation-background}

\paragraph{Multi-agent workflows.}
We represent a workflow as a directed acyclic graph $G=(V,E)$, whose nodes
are LLM requests and whose edges are data dependencies. An agent branch may
include multiple requests and tool calls; a successor runs only after its
dependencies are satisfied. WCT spans the release of the first request to
the completion of the last. For a given execution schedule, a request's
\emph{slack} is the maximum delay to its completion that does not postpone
workflow completion. Requests on a critical path have zero slack.

\paragraph{Serving and failure model.}
We consider a cluster of GPU workers with block-managed KV that can be
replicated to other GPUs or a host DRAM pool, consistent with KV-centric
architectures such as Mooncake~\cite{qin2025mooncake}. Our failure model is
a fail-stop failure of one GPU or worker: its local KV becomes unavailable,
while the control plane, backup DRAM, and other workers survive. Surviving
workers with the same model already loaded take over affected requests.
Whole-host failures, network partitions, and silent corruption are outside
this failure model.

\paragraph{Recovery paths and resource costs.}
Without a usable copy, the system must rerun prefill to rebuild lost KV.
A DRAM copy avoids reconstructing the saved prefix but must be transferred
to a GPU before use; concurrent recoveries share transfer bandwidth.
A copy already resident in the actual takeover GPU's HBM can be used
directly, at the cost of HBM capacity that could serve foreground requests
and replication traffic before failure. Backup placement therefore trades
capacity and normal-service overhead against recovery speed.

Prior systems reduce failure costs through KV streaming in
D\'{e}j\`{a}Vu~\cite{strati2024dejavu}, erasure coding in
GhostServe~\cite{jayakody2026ghostserve}, and load-aware checkpoint and
recovery coordination in LUMEN~\cite{cao2026lumen}. Multi-agent
applications add an allocation problem: saving the same recovery time for
different requests can have different value to the user. Neither saved KV
bytes nor avoided computation alone captures that difference.

\subsection{Recovery Cost Is Not Workflow Delay}
\label{sec:motivation-workflow}

A join waits for its last required branch. Earlier branches can tolerate
some delay, whereas delaying the branch that determines the join time
postpones downstream work. To characterize branch timing differences, we
analyze the public CC/Weka coding-agent
traces~\cite{semianalysis2026cctraces}. Of 393 sessions, 124 workflows
contain strictly overlapping subagent fan-outs: the latest branch start
precedes the earliest branch finish. For each fan-out, we compute the
longest minus the shortest branch duration; each workflow contributes its
maximum such gap. Figure~\ref{fig:motivation}(\subref{fig:motivation-branch-gap}) shows a median of
141\,s and a 90th percentile of 1,570\,s. These gaps establish substantial
branch-duration heterogeneity. With unequal branch start times, however,
a duration gap is neither join slack nor a direct measure of fault tolerance.

Figure~\ref{fig:motivation}(\subref{fig:motivation-workflow}) shows a short
five-branch fan-out, selected by minimizing the longest branch duration
among strictly overlapping five-way fan-outs. The branches take
40.6--59.1\,s; labels give completion times since workflow start.
A5 determines the join at 111.9\,s. Each branch's slack is its time
margin to this join.

Figure~\ref{fig:motivation}(\subref{fig:motivation-recovery-exposure}) shows
five independent counterfactuals at 76.0\,s after workflow start, each
losing one branch's current prefix KV. Trace hash-block counts give
prefix-length estimates of 57,984--95,744 tokens. We obtain recomputation
time $T_i$ by log--log interpolation of the measured MiniCPM4-8B curve in
Figure~\ref{fig:motivation-recompute} and derive branch slack $\sigma_i$
offline from trace completion times. Holding other branches, downstream
execution, and queueing unchanged, and counting only prefix recomputation,
the WCT penalty is $\max(0,T_i-\sigma_i)$. Substituting these costs
illustrates the effect of dependencies; it does not replay failures on the
original serving model.

A1 has the longest prefix and takes 15.4\,s to rebuild, but its 13.9\,s
slack limits the WCT penalty to 1.6\,s. A5 has the shortest prefix and
takes only 8.8\,s to rebuild, yet its zero slack exposes all of that time
to WCT. A2's 22.8\,s slack absorbs its entire 10.0\,s rebuild. Thus, avoiding the most
recomputation need not yield the largest WCT improvement. KV protection
must account for delay tolerance within a workflow; online estimation and
resource contention require the mechanisms and evaluation developed later.

\subsection{Long Contexts Increase Rebuilding Cost}
\label{sec:motivation-long-context}

\input{sections/motivation-prefill-figure}

Repeated agent interactions accumulate code, tool outputs, and other
agents' intermediate results in the context. A worker failure can force
this history through prefill again. Native sparse attention reduces
attention computation and access, but still needs the historical KV
required by subsequent execution; it does not eliminate state-loss costs.

Figure~\ref{fig:motivation-recompute} reports full-prefill costs for
8K--128K-token prefixes with no usable KV, where K~=~1,024. We run the
native sparse prefill path of MiniCPM4-8B~\cite{minicpm4model} on one
A100 80GB PCIe GPU with batch size one and BF16. Each point is the median
of five timings on the same random-token input, including block selection
and all KV writes. Measured latency grows from 1.05\,s at 8K to 22.3\,s
at 128K. For context, AIConfigurator~\cite{nvidia2026aiconfigurator}
0.9.0 HYBRID estimates at 128K are 34.6\,s for Qwen3-8B and 43.2\,s for
Qwen3-30B-A3B. These dashed curves use A100-SXM kernel tables and a
performance model; together with the PCIe measurements, they do not
constitute a controlled comparison of models or attention mechanisms.

Under these conditions, rebuilding a long context takes tens of seconds
even with native sparse attention. Combined with
Figure~\ref{fig:motivation}(\subref{fig:motivation-recovery-exposure}), this shows why both
rebuilding cost and branch delay tolerance matter. Saving KV before
failure can reduce repeated computation. The additional value of DRAM
or peer HBM depends on transfer, computation, and contention during
recovery and cannot be inferred from full-prefill latency alone.

\subsection{Problem and Design Goals}
\label{sec:motivation-goals}

AegisKV asks: \emph{with limited backup capacity and replication and
recovery bandwidth, which KV from which requests should receive what
protection to reduce the workflow-level impact of failures?}
Dependencies differentiate protection benefits, long contexts increase
state-loss costs, and replication competes with foreground serving.

This problem imposes five goals and constraints.
\textbf{(R1)} Optimize recovery for WCT, accounting for dependency-induced
delay tolerance.
\textbf{(R2)} Allocate limited backup capacity across requests and their
KV to reduce exposed recovery waiting.
\textbf{(R3)} Support concurrent recovery and transfer--compute overlap
without introducing request-level serialization for protection priority.
\textbf{(R4)} Bound normal-service replication overhead and subordinate
backup capacity and bandwidth use to foreground needs.
\textbf{(R5)} Preserve model computation and KV recovery correctness using
existing serving resources.
The next section develops the evidence for mechanism choices and the
corresponding placement and recovery design.
